\documentclass[oneside,a4paper,14pt]{extarticle}
\usepackage[T1,T2A]{fontenc}
\usepackage[a4paper,letterpaper,top=20mm,bottom=20mm,left=30mm,right=10mm]{geometry}
\usepackage[utf8]{inputenc}
\usepackage[russian]{babel}
\usepackage{textcomp}
\usepackage{indentfirst}
\usepackage{graphicx}
\usepackage{mwe}
\usepackage{wrapfig}
\usepackage{caption}
\usepackage{amsmath}
\usepackage{amsfonts}
\usepackage{amsthm}
\usepackage[all]{xy}
\usepackage[breaklinks]{hyperref}
\usepackage{titlesec}
\titleformat{\section}{\normalsize\bfseries}{\thesection}{1em}{}
\titleformat{\subsection}{\normalsize\bfseries}{\thesubsection}{1em}{}
\titleformat{\subsubsection}{\normalsize\bfseries}{\thesubsection}{1em}{}
\renewcommand\baselinestretch{1.33}\normalsize
\setlength{\parindent}{1.25cm}
\usepackage{indentfirst}

\begin{document}
    \newpage\thispagestyle{empty}
    \begin{center}
        МИНИСТЕРСТВО НАУКИ И ВЫСШЕГО ОБРАЗОВАНИЯ\\
            РОССИЙСКОЙ ФЕДЕРАЦИИ
            ФЕДЕРАЛЬНОЕ ГОСУДАРСТВЕННОЕ БЮДЖЕТНОЕ\\
            ОБРАЗОВАТЕЛЬНОЕ
            УЧРЕЖДЕНИЕ ВЫСШЕГО ОБРАЗОВАНИЯ\\
            «ВЯТСКИЙ ГОСУДАРСТВЕННЫЙ УНИВЕРСИТЕТ»\\
            Институт математики и информационных систем\\
            Факультет автоматики и вычислительной техники\\
            Кафедра электронных вычислительных машин
    \end{center}
    \vspace{20mm}

    \begin{center}
        Отчёт по лабораторной работе №4.2\\
        по дисциплине\\
        <<Программирование>>\\
    \end{center}
    \vspace{48mm}

    Выполнил студент гр. ИВТб-1302-06-00 \hspace{10mm} \rule[-0,5mm]{30mm}{0.15mm}\,/Изместьев М.Н./


    Проверил доцент кафедры ЭВМ \hfill  \rule[-0,5mm]{30mm}{0.15mm}\,/Баташев П.А./

    \vfill
    \begin{center}
        Киров\\
        2026
    \end{center}

    \newpage\thispagestyle{empty}

\begin{flushleft}
\section*{Цель работы:}
\noindent Научиться создавать простые графические интерфейсы с использованием библиотеки tkinter, организовывать обработку событий и интегрировать ранее разработанные программные модули в единое приложение. Закрепить понимание работы выбранной структуры данных через её визуальное представление.

\section*{Задание.}

1. На основе реализации структуры данных из лабораторной работы 4.1. необходимо разработать оконное приложение с графическим интерфейсом на Python.

2. Все операции со структурой данных должны выполняться исключительно через вызовы функций отдельно оформленных библиотечных модулей.

3. Необходимо предусмотреть реализацию, как минимум двух модулей: один на С++, второй --- Python.

4. Для получения дополнительных баллов, необходима реализация двух вариантов С++ модуля: с использованием структур на динамической памяти и с использованием контейнеров и адаптеров STL.

5. Интерфейс должен адекватно реагировать на попытку выполнить некорректную операцию (например, Pop из пустого стека), выводя понятное сообщение, без аварийного завершения программы.

6. Визуализация должна обновляться после каждой успешной операции.

7. В качестве ответа необходимо прикрепить отчёт и ссылку на репозиторий с проектом.

\textbf{Структура:} Кольцевая очередь (Circular Queue)

\textbf{Дополнительная возможность:} реализованы два варианта C++ модуля --- на динамической памяти и с использованием STL (\texttt{std::list})

\textbf{Язык:} C++, Python
\end{flushleft}

\newpage
\section*{Описание работы программы}

Приложение реализует графический интерфейс для работы с кольцевой очередью. Код разбит на три независимых модуля: Python-реализация, C++ на динамической памяти и C++ через STL. GUI позволяет переключаться между бэкендами в реальном времени.

\textbf{Принцип кольцевой очереди.} Структура хранит хвостовой указатель \texttt{tail}. Голова получается как \texttt{tail->next}. При добавлении новый узел вставляется после хвоста и становится новым хвостом; при удалении (dequeue) извлекается \texttt{tail->next}.

\textbf{GUI.} Написан в объектно-ориентированном стиле. Базовый интерфейс \texttt{QueueBackend} унифицирует работу с тремя реализациями. Класс \texttt{App} строит окно, обрабатывает события и перерисовывает \texttt{Canvas} после каждой операции. Элементы отображаются в виде цветных кружков на окружности со стрелками: зелёный --- HEAD, красный --- TAIL.

\newpage
\textbf{Модуль Python (circular\_queue\_py.py)}
\small
\begin{verbatim}
class Uzl:
    def __init__(self, data):
        self.data = data
        self.next = None


class Sruct:
    def __init__(self):
        self.tail = None
        self.size = 0

    def add(self, value):
        new_uzl = Uzl(value)
        if self.tail is None:
            new_uzl.next = new_uzl
            self.tail = new_uzl
        else:
            new_uzl.next = self.tail.next
            self.tail.next = new_uzl
            self.tail = new_uzl
        self.size += 1
        return 1

    def dequeue(self):
        if self.tail is None:
            return None
        head = self.tail.next
        value = head.data
        if self.tail == head:
            self.tail = None
        else:
            self.tail.next = head.next
        self.size -= 1
        return value

    def peek(self):
        if self.tail is None:
            return None
        return self.tail.next.data

    def clear(self):
        while self.size > 0:
            self.dequeue()

    def get_size(self):
        return self.size

    def get_elements(self):
        if self.tail is None:
            return []
        elements = []
        current = self.tail.next
        for _ in range(self.size):
            elements.append(current.data)
            current = current.next
        return elements
\end{verbatim}

\newpage
\textbf{Модуль C++ Dynamic (circular\_queue\_dynamic.cpp)}
\small
\begin{verbatim}
#ifdef _WIN32
#define EXPORT __declspec(dllexport)
#else
#define EXPORT
#endif

struct Uzl {
    int data;
    Uzl* next;
};

struct Sruct {
    Uzl* tail;
    int size;
};

static Sruct queue;

extern "C" {

EXPORT void lib_init() {
    queue.tail = nullptr;
    queue.size = 0;
}

EXPORT int lib_add(int value) {
    Uzl* newUzl = new Uzl;
    newUzl->data = value;
    if (queue.tail == nullptr) {
        newUzl->next = newUzl;
        queue.tail = newUzl;
    } else {
        newUzl->next = queue.tail->next;
        queue.tail->next = newUzl;
        queue.tail = newUzl;
    }
    queue.size++;
    return 1;
}

EXPORT int lib_dequeue(int* value) {
    if (queue.tail == nullptr) return 0;
    Uzl* head = queue.tail->next;
    *value = head->data;
    if (queue.tail == head) {
        queue.tail = nullptr;
    } else {
        queue.tail->next = head->next;
    }
    delete head;
    queue.size--;
    return 1;
}

EXPORT int lib_peek(int* value) {
    if (queue.tail == nullptr) return 0;
    *value = queue.tail->next->data;
    return 1;
}

EXPORT void lib_clear() {
    int value;
    while (queue.size > 0) lib_dequeue(&value);
}

EXPORT int lib_size() { return queue.size; }

EXPORT int lib_get_elements(int* buffer, int maxSize) {
    if (queue.tail == nullptr || maxSize <= 0) return 0;
    Uzl* current = queue.tail->next;
    int count = 0;
    for (int i = 0; i < queue.size && i < maxSize; i++) {
        buffer[i] = current->data;
        current = current->next;
        count++;
    }
    return count;
}

}
\end{verbatim}

\newpage
\textbf{Модуль C++ STL (circular\_queue\_stl.cpp)}
\small
\begin{verbatim}
#ifdef _WIN32
#define EXPORT __declspec(dllexport)
#else
#define EXPORT
#endif

#include <list>
static std::list<int> queue;

extern "C" {

EXPORT void lib_init()          { queue.clear(); }
EXPORT int  lib_add(int value)  { queue.push_back(value); return 1; }

EXPORT int lib_dequeue(int* value) {
    if (queue.empty()) return 0;
    *value = queue.front();
    queue.pop_front();
    return 1;
}

EXPORT int lib_peek(int* value) {
    if (queue.empty()) return 0;
    *value = queue.front();
    return 1;
}

EXPORT void lib_clear()  { queue.clear(); }
EXPORT int  lib_size()   { return (int)queue.size(); }

EXPORT int lib_get_elements(int* buffer, int maxSize) {
    if (queue.empty() || maxSize <= 0) return 0;
    int count = 0;
    for (auto it = queue.begin(); it != queue.end() && count < maxSize; ++it)
        buffer[count++] = *it;
    return count;
}

}
\end{verbatim}

\newpage
\textbf{GUI-приложение (gui\_app.py --- ключевые части)}
\small
\begin{verbatim}
import tkinter as tk
from tkinter import messagebox
import ctypes, os, sys, math, random

try:
    from circular_queue_py import Sruct
except ImportError:
    Sruct = None

class QueueBackend:
    def init(self): pass
    def add(self, val): pass
    def dequeue(self): return None
    def peek(self): return None
    def clear(self): pass
    def size(self): return 0
    def get_elements(self): return []

class PythonBackend(QueueBackend):
    def __init__(self):
        self.q = None
        self.is_list = Sruct is None
    def init(self):
        self.q = Sruct() if Sruct else []
    def add(self, val):
        if self.is_list: self.q.append(val)
        else: self.q.add(val)
    def dequeue(self):
        if self.is_list: return self.q.pop(0) if self.q else None
        return self.q.dequeue()
    def peek(self):
        if self.is_list: return self.q[0] if self.q else None
        return self.q.peek()
    def clear(self):
        if self.is_list: self.q.clear()
        else: self.q.clear()
    def size(self):
        return len(self.q) if self.is_list else self.q.get_size()
    def get_elements(self):
        return list(self.q) if self.is_list else self.q.get_elements()

class CppBackend(QueueBackend):
    def __init__(self, path):
        self.lib = None
        if os.path.exists(path):
            self.lib = ctypes.CDLL(path)
            self.lib.lib_init.restype      = None
            self.lib.lib_add.restype       = ctypes.c_int
            self.lib.lib_add.argtypes      = [ctypes.c_int]
            self.lib.lib_dequeue.restype   = ctypes.c_int
            self.lib.lib_dequeue.argtypes  = [ctypes.POINTER(ctypes.c_int)]
            self.lib.lib_peek.restype      = ctypes.c_int
            self.lib.lib_peek.argtypes     = [ctypes.POINTER(ctypes.c_int)]
            self.lib.lib_clear.restype     = None
            self.lib.lib_size.restype      = ctypes.c_int
            self.lib.lib_get_elements.restype  = ctypes.c_int
            self.lib.lib_get_elements.argtypes = [ctypes.POINTER(ctypes.c_int),
                                                   ctypes.c_int]
    def is_valid(self): return self.lib is not None
    def init(self):
        if self.lib: self.lib.lib_init()
    def add(self, val):
        if self.lib: self.lib.lib_add(val)
    def dequeue(self):
        if not self.lib: return None
        val = ctypes.c_int()
        res = self.lib.lib_dequeue(ctypes.byref(val))
        return val.value if res else None
    def peek(self):
        if not self.lib: return None
        val = ctypes.c_int()
        res = self.lib.lib_peek(ctypes.byref(val))
        return val.value if res else None
    def clear(self):
        if self.lib: self.lib.lib_clear()
    def size(self):
        return self.lib.lib_size() if self.lib else 0
    def get_elements(self):
        if not self.lib: return []
        buf = (ctypes.c_int * 10000)()
        count = self.lib.lib_get_elements(buf, 10000)
        return [buf[i] for i in range(count)]
\end{verbatim}

\newpage
\small
\begin{verbatim}
class App:
    def __init__(self, root):
        self.root = root
        self.backend = None
        self.is_created = False
        base = os.path.dirname(os.path.abspath(__file__))
        ext = ".dll" if sys.platform == "win32" else ".so"
        self.backends = {
            "Python":     PythonBackend(),
            "C++ Dynamic": CppBackend(os.path.join(base,
                              "circular_queue_dynamic" + ext)),
            "C++ STL":     CppBackend(os.path.join(base,
                              "circular_queue_stl" + ext))
        }
        self.setup_ui()

    def setup_ui(self):
        top = tk.Frame(self.root)
        top.pack(fill=tk.X, padx=10, pady=5)
        tk.Label(top, text="Бэкенд:").pack(side=tk.LEFT)
        self.bk_var = tk.StringVar(value="Python")
        for name, bk in self.backends.items():
            if name == "Python" or (isinstance(bk, CppBackend) and bk.is_valid()):
                tk.Radiobutton(top, text=name, variable=self.bk_var,
                    value=name, command=self.on_change_backend).pack(side=tk.LEFT)
        self.lbl_status = tk.Label(self.root, text="Не создано", fg="red")
        self.lbl_status.pack()
        self.canvas = tk.Canvas(self.root, width=880, height=380,
                                bg="white", relief=tk.SUNKEN, bd=2)
        self.canvas.pack(padx=10, pady=5)
        inp = tk.Frame(self.root)
        inp.pack(fill=tk.X, padx=10)
        self.ent_val = self.mk_entry(inp, "Значение:", 8)
        self.ent_cnt = self.mk_entry(inp, "Кол-во:",   5)
        self.ent_min = self.mk_entry(inp, "Мин:",      5)
        self.ent_max = self.mk_entry(inp, "Макс:",     5)
        btns = tk.Frame(self.root)
        btns.pack(fill=tk.X, padx=10, pady=5)
        self.mk_btn(btns, "Создать",   self.create)
        self.mk_btn(btns, "Добавить",  self.add)
        self.mk_btn(btns, "Удалить",   self.dequeue)
        self.mk_btn(btns, "Peek",      self.peek)
        self.mk_btn(btns, "Случайные", self.rand_add)
        self.mk_btn(btns, "Очистить",  self.clear)
        self.log_txt = tk.Text(self.root, height=6,
                               state=tk.DISABLED, bg="#f0f0f0")
        self.log_txt.pack(fill=tk.BOTH, padx=10, pady=5)

    def refresh(self):
        self.canvas.delete("all")
        if not self.is_created:
            self.draw_text("Структура не создана")
            return
        elems = self.backend.get_elements()
        sz = len(elems)
        self.lbl_status.config(
            text="Бэкенд: " + self.bk_var.get() + " | Размер: " + str(sz),
            fg="green")
        if sz == 0:
            self.draw_text("Очередь пуста")
            return
        cx, cy = 440, 200
        r  = 30 + sz * 12 if sz < 10 else 150
        nr = 22
        coords = []
        for i in range(sz):
            ang = -math.pi / 2 + 2 * math.pi * i / sz
            coords.append((cx + r * math.cos(ang), cy + r * math.sin(ang)))
        for i in range(sz):
            x1, y1 = coords[i]
            x2, y2 = coords[(i + 1) % sz]
            dx, dy = x2 - x1, y2 - y1
            dist = math.sqrt(dx*dx + dy*dy)
            if dist == 0: continue
            self.canvas.create_line(
                x1 + dx/dist*nr, y1 + dy/dist*nr,
                x2 - dx/dist*nr, y2 - dy/dist*nr,
                arrow=tk.LAST, fill="#555", width=2)
        for i, (x, y) in enumerate(coords):
            col = "#4CAF50" if i == 0 else ("#F44336" if i == sz-1 else "#2196F3")
            self.canvas.create_oval(x-nr, y-nr, x+nr, y+nr, fill=col, width=2)
            self.canvas.create_text(x, y, text=str(elems[i]),
                                    fill="white", font=("Arial", 10, "bold"))
        self.canvas.create_text(coords[0][0], coords[0][1]-nr-12,
                                text="HEAD", fill="green",
                                font=("Arial", 9, "bold"))
        self.canvas.create_text(coords[-1][0], coords[-1][1]+nr+12,
                                text="TAIL", fill="red",
                                font=("Arial", 9, "bold"))

if __name__ == "__main__":
    root = tk.Tk()
    root.title("Циклическая очередь")
    root.geometry("900x720")
    App(root)
    root.mainloop()
\end{verbatim}

\normalsize
\newpage
\textbf{Экранные формы:}

\centering
    \includegraphics[width=0.85\linewidth]{1.png}
    \label{fig:screen1}

    \textbf{Рисунок 1 --- }Главное окно, Python-бэкенд, пустая очередь

\centering
    \includegraphics[width=0.85\linewidth]{2.png}
    \label{fig:screen2}

    \textbf{Рисунок 2 --- }Элементы в очереди (визуализация на Canvas)

\centering
    \includegraphics[width=0.85\linewidth]{3.png}
    \label{fig:screen3}

    \textbf{Рисунок 3 --- }Переключение на C++ Dynamic

\centering
    \includegraphics[width=0.85\linewidth]{4.png}
    \label{fig:screen4}

    \textbf{Рисунок 4 --- }Операция Peek

\centering
    \includegraphics[width=0.85\linewidth]{5.png}
    \label{fig:screen5}

    \textbf{Рисунок 5 --- }Сообщение об ошибке при операции с пустой очередью

\begin{flushleft}
\newpage
\section*{Вывод}
\noindent
В ходе лабораторной работы я научился создавать графический интерфейс с использованием библиотеки tkinter и интегрировать в него модули на Python и C++. Была реализована кольцевая очередь в трёх вариантах: Python, C++ на динамической памяти и C++ через STL (\texttt{std::list}). GUI позволяет переключаться между бэкендами в реальном времени и визуализирует состояние структуры на Canvas в виде кольца с направленными стрелками. Все некорректные операции обрабатываются без аварийного завершения программы.
\end{flushleft}

\end{document}
